% English translation of checked_pages/esq_119_034.tex; source: cote 119, site 34, PDF page 35, typescript page 32.
% AI-translated by GPT-6 Astra (OpenAI, ChatGPT Pro mode).
\scanpage{119/34}
\mspage{32}

independent of the tame theory chosen to express
it.

The first decisive test for a good system of axioms
for a notion of a ``tame subset of $\mathbb{R}^n$ '' seems to me to be the possibility
of proving such comparison theorems. So far I have
been content to draw out a plausible provisional system of axioms,
without any assurance that further axioms will not have to be
added, which probably only ``work on concrete examples''\tnote{TN-034-01}{``Work on concrete examples'' renders ``travail sur pièces'', the author's phrase for work that exposes the axioms needed in actual instances.}
will bring to light. The strongest of the axioms I have
introduced, and probably the one whose verification in specific
cases is (or will be) most delicate, is an \uline{axiom of}
\uline{triangulability} (tame, it goes without saying) of a tame subset of $\mathbb{R}^n$.
I have not attempted to prove the comparison theorem from these axioms
alone; nevertheless I have had the impression (rightly
or wrongly again !) that this proof, whether or not it requires
the introduction of some additional axiom, will present
no great technical difficulty. It is quite possible
that the technical difficulties in developing
satisfactory foundations for tame topology, including a
theory of dévissage of stratified tame structures, are
already concentrated for the most part in the axioms, and consequently
have essentially already been overcome by triangulability
theorems in the manner of Lojasiewicz and Hironaka. What is lacking,
once again, is by no means mathematicians' sometimes impressive
technical virtuosity, but the daring (or simply\frapparatus{V-FR-034-001}
the innocence ...) to free oneself from a familiar context
accepted by an unbroken consensus ...

The advantages of an axiomatic approach toward foundations
for tame topology seem fairly obvious to me. Thus,
to regard a complex algebraic variety, or the set of
real points of an algebraic variety defined over $\mathbb{R}$, as a
tame space, it seems preferable to work in the
``piecewise $\mathbb{R}$-algebraic'' theory, or even in the piecewise $\overline{\mathbb{Q}}_r$-algebraic
theory (where $\overline{\mathbb{Q}}_r = \overline{\mathbb{Q}} \cap \mathbb{R}$) when dealing with varieties defined

\margin{34}

